\section{Quasi-Newton BFGS}

\subsection{General explanation}

The initial idea of the BFGS algorithm was to accelerate the iteration process in a nonlinear optimisation because the computer material wasn’t stable. The improvement developed by this method is based on the fact that “quasi-Newton methods […] require only the gradient of the objective function to be supplied at each iterate” \cite{nocedalwright1999}. This method is named for its discoverers called Broyden, Fletcher, Goldfarb, and Shanno \cite{nocedalwright1999}.

Unlike Newton’s method, which required explicit computation of the Hessian matrix, BFGS builds an approximation of the Hessian using only gradient information collected during the optimisation process. This method generates a sequence of iterates $x_k$ starting from an initial guess $x_0$. At each iteration, it computes a search direction $p_k$ using the current Hessian approximation $B_k$ and gradient $\nabla f_k$, then performs a line search to determine the step length $\alpha_k$. The algorithm terminates when convergence criteria are satisfied. This happens typically when the gradient norm falls below a specified tolerance \cite{nocedalwright1999}.

\subsection{Mathematical approach}

The BFGS method approximates the inverse Hessian $H_k \approx (\nabla^2 f_k)^{-1}$ and updates it at each iteration using the following formula. It was independently derived by Fletcher \cite{fletcher1970} and demonstrated by Shanno \cite{shanno1970}.

\begin{equation}
H_{k+1} = (I - \rho_k s_k y_k^T) H_k (I - \rho_k y_k s_k^T) + \rho_k s_k s_k^T
\end{equation}

where:
\begin{align*}
s_k &= x_{k+1} - x_k \quad \text{(step taken)} \\
y_k &= \nabla f_{k+1} - \nabla f_k \quad \text{(change in gradient)} \\
\rho_k &= \frac{1}{y_k^T s_k}
\end{align*}

The search direction is computed as: 
\begin{equation}
p_k = -H_k \nabla f_k
\end{equation}

The new iterate is obtained via line search: 
\begin{equation}
x_{k+1} = x_k + \alpha_k p_k
\end{equation}
where $\alpha_k$ satisfies the Wolfe conditions to ensure sufficient decrease and curvature conditions.

For the Two-Bar Truss problem with constraints, the penalized objective function becomes:
\begin{equation}
F(x) = w_f(f_1(x) + f_2(x)) + w_c \times (\text{penalty}_{g1} + \text{penalty}_{g2})
\end{equation}
where $\text{penalty}_{gi} = \max(0, g_i(x))^2$ for $i = 1, 2$.
